% Page 022

\seite{22}

konnte man Heu- und Körnerernte um das Doppelte übertreffen.\ob
Der Kartoffelertrag war groß und die Qualität gut.\ob
Obst war verhältnismäßig viel. In Betracht kommen\ob
die wenigen Obstbäume und die geographische Lage.\ob
Durch die Menge des gewachsenen Obstes\unsicher{} wurde der\ob
Preis so herunter gedrückt, daß daraus kaum die\ob
Arbeit bezahlt wurde.

\pend
\abschnitt{Schulfeiern}
\pstart

Das Schuljahr 1895/96 war reich an Schulfeiern. Am\ob
2.\ 9. 1895 war die 25jährige Wiederkehr der glorreichen\ob
Schlacht bei Sedan; am 18.\ Januar 1896 das 25jährige\ob
Bestehen des neuerrichteten Deutschen Kaiserreiches.\ob
Am 27.\ Januar war der Geburtstag Sr. Majestät\ob
Kaiser Wilhelms II. An genannten Tagen fand eine\ob
dem Feste entsprechende Feier statt. Am Geburtstage\ob
Seiner Majestät wurde die Feier im Schulsaale zu\ob
Seffern abgehalten.

Am 1.\ April, am achtzigsten Geburtstag des Fürsten\ob
Bismarck, fiel zur Feier des Tages der Schulunter\ohb
richt aus.

Im Jahre 1895 wurden 2 Volkszählungen vorgenommen.\ob
Die erste, im Juni, war eine Gewerbezählung. Die\ob
zweite, am 2.\ Dezember war die eigentliche Volkszählung.\ob
Bei dieser hat das Dorf Sefferweich gegenwärtig 50 Häuser.\ob
Diese sind dem landwirthschaftlichen Betriebe seiner Bewohner
\pend
